% !PRINTED 460
\weberpage{460}
%<<CONT

\[
\begin{aligned}
D &= -8, & \omega &= \sqrt{-2}, \\
D &= -7, & \omega &= \frac{-3+\sqrt{-7}}{2}, \\
D &= -4, & \omega &= i.
\end{aligned}
\]

The value $\gamma_2(i)$ is, however, known (by \S\,117), namely
$=12$, and therefore the left side of (14) must be divisible by
$u-12$. Since this factor is known, the others are easily found:

\[
u^4-2u^3-495u^2+2^4\cdot 3^3\cdot 5^3 = (u-12)(u-20)(u+15)^2.
\]

Since $\gamma_2(\omega)$, just like $j(\omega)$, has a positive
value for a purely imaginary $\omega$, the factor $u-20$ must
correspond to the value $D = -8$, and we obtain:

\weq{16}{\gamma_2(i) = 12,}

\weq{17}{\gamma_2(\sqrt{-2}) = 20,}

\weq{18}{\gamma_2\left(\frac{-3+\sqrt{-7}}{2}\right) = -15.}

Besides $-7$ there exist five further odd negative discriminants
not divisible by 3:

\weq{19}{-11, \quad -19, \quad -43, \quad -67, \quad -163^{1)},}

for which there is only \gesperrt{one} class of forms. For these, then, by
our theorem:

\[
-\gamma_2\left(\frac{-3+\sqrt{-m}}{2}\right)
\]

is an \gesperrt{integral rational number} $Z$.

In order to find these rational numbers, we shall seek limits
between which they must lie, and which differ from one another by
less than a unit.

By \S\,54, (5), (8) we have:

\[
\gamma_2(\omega) = f(\omega)^{16}-\frac{16}{f(\omega)^8} = f_1(\omega)^{16}+\frac{16}{f_1(\omega)^8} = f_2(\omega)^{16}+\frac{16}{f_2(\omega)^8},
\]

further, by \S\,34, (19):

\[
f_2\left(\frac{-3+\omega}{2}\right) f(\omega) = e^{-\frac{\pi i}{8}}\sqrt{2},
\]

\bigskip
\noindent\rule{4cm}{0.4pt}

\footnotesize
1) That no further discriminants of this kind exist can so far be
inferred only from the fact that, as far as the computation of the
class numbers has been carried, no further numbers of this kind
have been found [cf. the table of class numbers of
\gesperrt{Gauss} (Werke, Bd.\,II, S.\,450)].
\normalsize

%CONT>>
